% !TeX root = ../../Book1.tex
% 纲要标题：### 12.1 向量丛及其粘合
% 纲要位置：几何学计划纲要.md，第 572 行。
% 正文状态：初稿；正文、证明、例子与配套习题解答已撰写，尚未正式审核。

% 纲要内容开始
% 
% * 局部平凡化、过渡函数和余循环
% * 丛同态、子丛及常秩条件
% * 拉回丛、直和与商丛
% * 标架丛与全局标架
% 
% 纲要内容结束

\section{向量丛及其粘合}
\label{g1:sec:1201}

切丛把每点的切空间组织成一个光滑对象。它的局部坐标同时记住基点与向量分量，而换坐标时，分量按随基点光滑变化的可逆矩阵变换。许多几何对象也有这种形式：曲面每点的法向量、线性方程在每点的解空间，以及随后要讨论的张量，都不是固定向量空间中的一组数，而是一族随基点变化的向量空间。向量丛保留了切丛构造中的这一共同部分。

\ProseParagraphBreak

本节讨论有限秩实光滑向量丛，底空间为 Hausdorff、第二可数的光滑流形，允许有边界；边界附近的光滑性沿用局部延拓到 Euclid 开集的定义。我们先弄清局部数据怎样粘合，再讨论如何逐纤维取子空间、商空间，以及如何更换底空间。关键问题始终是：逐点成立的线性代数，何时能随基点光滑地进行。

\subsection{局部平凡化、过渡函数和余循环}
\label{g1:subsec:1201:01}

\cref{g1:thm:tangent-bundle-atlas}已经证明，切丛在坐标域 \(U\) 上可写成 \(U\times\mathbb R^n\)，并验证了总空间的 Hausdorff 性与第二可数性。这种乘积表示只在基点的一个邻域上给定；换图时可以出现非恒定矩阵，因此它还没有给出整个切丛与 \(M\times\mathbb R^n\) 的同构。

\begin{definition}
	\label{g1:def:vector-bundle}
	设 \(M\) 为 Hausdorff、第二可数的光滑 \(n\) 维流形，可带边界，\(r\geq0\) 为整数。\(M\) 上的秩 \(r\) 的实\term[xiangliangcong]{向量丛}[vector bundle]由光滑流形 \(E\)、光滑满射 \(\pi:E\to M\) 及每个集合 \(E_p=\pi^{-1}(p)\) 上的 \(r\) 维实向量空间结构组成，并要求每个 \(p\in M\) 都有开邻域 \(U\) 和微分同胚
	\[
		\phi_U:\pi^{-1}(U)\longrightarrow U\times\mathbb R^r,
		\qquad \phi_U(v_p)=(p,[v_p]_U),
	\]
	使 \([\cdot]_U:E_q\to\mathbb R^r\) 对每个 \(q\in U\) 都是线性同构。这样的 \(\phi_U\) 称为\term[jubupingfanhua]{局部平凡化}[local trivialization]。\(E\)、\(M\)、\(E_p\) 分别称为总空间、底空间和 \(p\) 上的\term[xianwei]{纤维}[fiber]；\(E|_U=\pi^{-1}(U)\) 表示限制丛。
\end{definition}

这里总空间同样要求 Hausdorff、第二可数。局部模型说明 \(\dim E=n+r\)；若底空间有边界，模型为 \(\mathbb H^n\times\mathbb R^r\)，非负坐标来自底空间，所以 \(\partial E=\pi^{-1}(\partial M)\)。秩零的丛每个纤维只有零向量，其总空间可由 \(\pi\) 识别为 \(M\)。本文固定一个秩；若底空间不连通，也可在不同连通分支分别规定秩。

\begin{definition}
	\label{g1:def:bundle-section}
	设 \(\pi:E\to M\) 为有限秩实光滑向量丛，\(U\subset M\) 开。满足 \(\pi\circ s=\operatorname{id}_U\) 的光滑映射 \(s:U\to E\) 称为 \(E\) 在 \(U\) 上的\term[jiemian]{截面}[section]，全体光滑截面记为 \(\Gamma(U,E)\)，并简记 \(\Gamma(E)=\Gamma(M,E)\)。若秩为 \(r\)，一组截面 \(e_1,\ldots,e_r\) 在每个 \(p\in U\) 都组成 \(E_p\) 的有序基，就称为 \(U\) 上的\term[jububiaojia]{局部标架}[local frame]。
\end{definition}

局部平凡化把标准基向量送回纤维，便给出局部标架。反过来，给定局部标架 \(e=(e_1,\ldots,e_r)\)，记
\[
	e(p)v=\sum_{a=1}^r v^a e_a(p),\qquad v\in\mathbb R^r.
\]
在任一已有平凡化中，\(e_a\) 的分量组成一个光滑可逆矩阵 \(A(p)\)，因而 \((p,v)\mapsto e(p)v\) 及其逆都光滑。逆的光滑性来自 \(A^{-1}=\operatorname{adj}(A)/\det A\)，这里行列式处处非零。于是局部标架与局部平凡化确实等价，而不仅是逐点选基。

\ProseParagraphBreak

在开覆盖 \(\{U_i\}\) 上选标架 \(e_i\)，其中下标 \(i\) 标记开集，标架内部的向量仍用 \(e_{i,a}\) 表示。在交集上存在唯一光滑函数
\[
	g_{ij}:U_i\cap U_j\longrightarrow\operatorname{GL}(r,\mathbb R)
\]
使
\begin{equation}
	\label{g1:eq:bundle-frame-transition}
	e_j=e_i g_{ij},\qquad v_i=g_{ij}v_j.
\end{equation}
第二式中 \(v_i,v_j\) 是同一个纤维向量在两组标架下的列分量，因为 \(e_i v_i=e_jv_j\)。因此标架变换与分量变换在这个下标约定下都出现 \(g_{ij}\)，但前一式把第 \(j\) 组基写成第 \(i\) 组基的线性组合，后一式则把第 \(j\) 组分量转换成第 \(i\) 组分量。

\ProseParagraphBreak

这些\term[guoduhanshu]{过渡函数}[transition function]满足
\begin{equation}
	\label{g1:eq:bundle-cocycle}
	g_{ii}=I_r,\qquad g_{ji}=g_{ij}^{-1},\qquad
	g_{ij}g_{jk}=g_{ik}.
\end{equation}
最后一个等式在三重交集上成立，称为\term[yuxunhuantiaojian]{余循环条件}[cocycle condition]：先由第 \(k\) 组分量换到第 \(j\) 组，再换到第 \(i\) 组，必须与直接转换相同。它由 \(e_k=e_jg_{jk}=e_i g_{ij}g_{jk}\) 及基的线性无关性得到。对截面 \(s\)，局部表达 \(s=e_i s_i\) 中的列函数 \(s_i:U_i\to\mathbb R^r\) 光滑，且满足 \(s_i=g_{ij}s_j\)；反之，满足这些关系的光滑列函数唯一确定一个光滑截面。纤维中的加法与数乘由此定义了截面的逐点运算，并使 \(\Gamma(E)\) 成为 \(C^\infty(M)\) 上的模。

\begin{example}
	\label{g1:exm:tangent-bundle-transition-functions}
	对光滑流形 \(M\)，在坐标图 \(x_i:U_i\to x_i(U_i)\) 上取切标架 \(e_i=(\partial_{x_i^1},\ldots,\partial_{x_i^n})\)。由\eqref{g1:eq:tangent-bundle-transition}，
	\[
		g_{ij}(p)=\mathrm{D}(x_i\circ x_j^{-1})(x_j(p)).
	\]
	链式法则正好给出 \(g_{ij}g_{jk}=g_{ik}\)，所以切丛的过渡函数是坐标变换的 Jacobian 矩阵。一般向量丛的矩阵只须光滑、可逆并满足余循环条件，无须来自某个底空间坐标变换的导数。乘积丛 \(M\times\mathbb R^r\) 则有处处相同的标准标架，过渡函数都可取为 \(I_r\)。
\end{example}

余循环条件足以把一批局部乘积真正粘成向量丛。证明中既要保证同一个向量不会因走不同的换图路径而得到不同代表，也要保证粘合后的空间满足流形的整体拓扑条件。

\begin{theorem}[向量丛的粘合]
	\label{g1:thm:bundle-cocycle-gluing}
	设 \(M\) 为 Hausdorff、第二可数的光滑流形，可带边界；\(\{U_i\}_{i\in I}\) 为开覆盖，\(r\geq0\) 为固定整数。若光滑映射 \(g_{ij}:U_i\cap U_j\to\operatorname{GL}(r,\mathbb R)\) 在相应的二重、三重交集上满足 \(g_{ii}=I_r\)、\(g_{ji}=g_{ij}^{-1}\)、\(g_{ij}g_{jk}=g_{ik}\)，即\eqref{g1:eq:bundle-cocycle}，则它们确定一个秩 \(r\) 的实光滑向量丛，其局部平凡化按 \((p,v_j)\mapsto(p,g_{ij}(p)v_j)\) 转换。该丛在保持这些局部识别的丛同构意义下唯一。
\end{theorem}
\begin{proof}
	先把各局部乘积视为互不相交的副本，令
	\[
		\mathcal E=\coprod_{i\in I}(U_i\times\mathbb R^r),\qquad
		(i,p,v)\sim(j,q,w)
		\quad\Longleftrightarrow\quad
		p=q\ \text{且}\ v=g_{ij}(p)w.
	\]
	自反性来自 \(g_{ii}=I_r\)，对称性来自 \(g_{ji}=g_{ij}^{-1}\)。若 \(v=g_{ij}(p)w\)、\(w=g_{jk}(p)z\)，则余循环条件给出 \(v=g_{ik}(p)z\)，所以关系也传递。令 \(E=\mathcal E/\!\sim\)，赋予商拓扑，并记商映射为 \(Q\)。由于等价关系不改变基点，\(\pi([i,p,v])=p\) 良定义且连续。

	商映射 \(Q\) 是开映射。事实上，对 \(\mathcal E\) 中的开集 \(O\)，它的饱和集在第 \(j\) 个副本中的部分为
	\[
		\bigcup_{i\in I}\left\{(p,w):p\in U_i\cap U_j,\ (p,g_{ij}(p)w)\in O_i\right\},
	\]
	其中 \(O_i=O\cap(U_i\times\mathbb R^r)\)。每项都是光滑同胚 \((p,w)\mapsto(p,g_{ij}(p)w)\) 下的开原像，故此并为开。商拓扑的定义于是说明 \(Q(O)\) 开。

	对每个 \(i\)，映射
	\[
		\tau_i:U_i\times\mathbb R^r\longrightarrow\pi^{-1}(U_i),
		\qquad\tau_i(p,v)=[i,p,v]
	\]
	是双射：每个基点在 \(U_i\) 的等价类都可换到第 \(i\) 个副本，而在同一副本中相等的两个代表必有相同分量。它连续且开，故为同胚；其像 \(\pi^{-1}(U_i)\) 也是 \(E\) 中的开集。在交集上有
	\[
		\tau_i^{-1}\tau_j(p,w)=(p,g_{ij}(p)w).
	\]
	再将每个 \(U_i\) 限制到底空间坐标域，便得到 \(E\) 的光滑坐标图；上述公式以及底空间换图的光滑性保证它们相容。在边界附近，矩阵元有局部光滑延拓，乘法公式同样给出半空间模型之间的光滑映射。

	还需核验 Hausdorff 性。两个点若有不同基点，用底空间中不交开邻域的逆像将它们分离；若基点相同，就把两点放进同一个 \(\tau_i\) 的像内，利用 \(U_i\times\mathbb R^r\) 的 Hausdorff 性分离。该像是 \(E\) 的开集，所以所得邻域也是总空间的开邻域。

	第二可数性不能由任意多个局部图直接推出，不过底空间的第二可数性提供可数子覆盖。具体地，对底空间可数基中每个包含在某个 \(U_i\) 内的基元素，选择一个包含它的 \(U_i\)；这些被选开的集合构成可数覆盖。每个对应乘积 \(U_i\times\mathbb R^r\) 有可数基，把它们通过 \(\tau_i\) 送到 \(E\) 再取并，就得到总空间的可数基。

	在纤维 \(E_p\) 中，选任一个 \(p\in U_i\)，用 \(\tau_i(p,\cdot)\) 运输 \(\mathbb R^r\) 的向量空间运算。换到另一副本时使用的是线性同构 \(g_{ij}(p)\)，故运算与所选副本无关。投影在局部是 \((p,v)\mapsto p\)，光滑且满射，\(\tau_i^{-1}\) 则是所需局部平凡化。

	最后，若另一个丛 \(E'\) 有同样局部识别 \(\tau_i'\)，则 \([i,p,v]\mapsto\tau_i'(p,v)\) 由相同过渡关系而良定义。它在每个局部乘积上为恒等坐标映射，故是光滑丛同构。任何保持局部识别的映射都必须使用这个公式，因此唯一。秩零时取 \(\operatorname{GL}(0,\mathbb R)=\{I_0\}\)，所有论证仍成立。
\end{proof}

这里的唯一性保留了指定的局部识别；改变标架当然会改变矩阵。若在 \(U_i\) 上以 \(\widetilde v_i=h_i v_i\) 改换分量，其中 \(h_i:U_i\to\operatorname{GL}(r,\mathbb R)\) 光滑，则
\begin{equation}
	\label{g1:eq:bundle-cocycle-change}
	\widetilde e_i=e_i h_i^{-1},\qquad
	\widetilde g_{ij}=h_i g_{ij}h_j^{-1}.
\end{equation}
后一个等式由 \(\widetilde v_i=h_i g_{ij}h_j^{-1}\widetilde v_j\) 得出。反过来，两组余循环若满足此式，\([i,p,v]\mapsto[i,p,h_i(p)v]\) 就给出所粘丛之间的同构：把 \(v=g_{ij}w\) 代入，得到 \(h_i v=\widetilde g_{ij}h_jw\)，恰好保证映射与代表选择无关。不同覆盖上的表示可先限制到两覆盖的交集，再作同一比较。因此过渡矩阵是丛在选定标架中的数据，并不是无需改换的绝对记号。

\begin{example}[Möbius 线丛]
	\label{g1:exm:mobius-line-bundle}
	令 \(M=S^1=\{e^{\mathrm{i}t}:t\in\mathbb R\}\)，把 \(\mathbb R\times\mathbb R\) 按
	\[
		(t+2\pi,v)\sim(t,-v)
	\]
	及其迭代关系取商，记所得空间为 \(L\)，投影为 \([t,v]\mapsto e^{\mathrm{i}t}\)。每个足够短的圆弧 \(U_i\) 上可选光滑角度提升 \(t_i:U_i\to\mathbb R\)，从而有局部表示 \((p,v)\mapsto[t_i(p),v]\)。若在交叠上 \(t_j-t_i=2\pi k_{ij}\)，则 \(k_{ij}\) 为局部常值整数，且
	\[
		[t_j(p),v]=[t_i(p),(-1)^{k_{ij}(p)}v].
	\]
	所以过渡函数为 \(g_{ij}=(-1)^{k_{ij}}\)，满足余循环条件。\cref{g1:thm:bundle-cocycle-gluing}给出秩一光滑丛结构；这些局部表示也是所写商空间的图，因为任一长度小于 \(2\pi\) 的提升圆弧及其 \(2\pi\) 平移副本，经商映射的饱和集仍开。这个丛称为 Möbius 线丛。

	它的光滑截面恰对应满足
	\begin{equation}
		\label{g1:eq:mobius-section-antiperiodic}
		f(t+2\pi)=-f(t)
	\end{equation}
	的光滑函数 \(f:\mathbb R\to\mathbb R\)，截面公式为 \(s(e^{\mathrm{i}t})=[t,f(t)]\)。因为 \([t+2\pi,f(t+2\pi)]=[t,-f(t+2\pi)]\)，此条件正是代表选择无关性；光滑性可在各提升圆弧上核验。若截面处处非零，\(f(0)\) 与 \(f(2\pi)=-f(0)\) 异号，中间值定理迫使某处 \(f\) 为零，矛盾。故 \(L\) 没有处处非零的光滑截面，甚至没有这样的连续截面。乘积线丛却有常值分量为 \(1\) 的截面，任何丛同构都会保持非零性，因此 \(L\) 与 \(S^1\times\mathbb R\) 不同构。
\end{example}

\begin{figure}[htbp]
	\centering
	% !TeX root = ../Book1.tex
% Möbius线丛的端点粘合；纤维只截取有限区间用于示意。
\begin{tikzpicture}[x=0.95cm,y=1.15cm,>=stealth]
	\fill[blue!3] (0,-1.35) rectangle (6.283,1.35);
	\draw[gray!65,dashed] (0,1.35)--(6.283,1.35) (0,-1.35)--(6.283,-1.35);
	\draw[very thick,->] (0,-1.35)--(0,1.35);
	\draw[very thick,->] (6.283,1.35)--(6.283,-1.35);
	\draw[gray!65,->] (-0.3,0)--(6.8,0) node[right] {\(t\)};
	\node[above] at (0,1.4) {\(v\)};
	\node[below left] at (0,0) {\(0\)};
	\node[below] at (3.142,-0.04) {\(\pi\)};
	\node[below right] at (6.283,0) {\(2\pi\)};
	\draw[red!75!black,dashed,thick] (0,0.55)--(6.283,0.55);
	\draw[red!75!black,fill=white] (0,0.55) circle (2pt) (6.283,0.55) circle (2pt);
	\node[red!75!black,above] at (4.9,0.6) {恒正分量不相合};
	\draw[blue!75!black,very thick,domain=0:6.283,samples=80] plot (\x,{cos(deg(\x/2))});
	\fill[blue!75!black] (0,1) circle (2pt) (3.142,0) circle (2pt) (6.283,-1) circle (2pt);
	\node[blue!75!black,above right] at (0.5,1.02) {\(f(t)=\cos(t/2)\)};
	\node[left] at (0,1) {\(1\)};
	\node[right] at (6.283,-1) {\(-1\)};
	\node[below,align=center] at (3.142,-1.52) {\((2\pi,v)\sim(0,-v)\)：两侧沿相反方向粘合};
\end{tikzpicture}

	\caption[Möbius 线丛中的截面与端点粘合]{矩形只画出线丛纤维的有限一段。左右竖边按相反方向粘合，\((2\pi,v)\) 对应 \((0,-v)\)。蓝线为截面分量 \(f(t)=\cos(t/2)\)，它的两端经过粘合相合，而在中间穿过零。红色虚线的分量恒为正，左右端点却不能按粘合关系相合，因而不是圆周上的截面。}
	\label{g1:fig:mobius-section}
\end{figure}

这说明局部平凡并未排除全局扭转。扭转并不是某个纤维内部的缺陷：每条纤维都是一维实向量空间；障碍出现在沿圆周比较各条纤维的基时。

\subsection{丛同态、子丛及常秩条件}
\label{g1:subsec:1201:02}

对同一个底空间上的两族向量空间，可以逐点施加线性映射。不过逐点线性只保证每个纤维内的行为，并不保证映射随基点连续或光滑。局部标架把后一个要求化为矩阵元的光滑性。

\begin{definition}
	\label{g1:def:bundle-morphism}
	设 \(E\to M\)、\(F\to M\) 为有限秩实光滑向量丛。覆盖恒等映射的\term[congtongtai]{丛同态}[bundle morphism]是光滑映射 \(\Phi:E\to F\)，使 \(\pi_F\Phi=\pi_E\)，且每个限制 \(\Phi_p:E_p\to F_p\) 都线性。如果它有同样性质的光滑逆，称为\term[congtonggou]{丛同构}[bundle isomorphism]。若 \(E\to N\)、\(F\to M\) 的底空间不同，覆盖光滑映射 \(f:N\to M\) 的丛同态则要求 \(\pi_F\Phi=f\pi_E\)，并要求 \(E_q\to F_{f(q)}\) 线性。
\end{definition}

\begin{proposition}[丛同态的矩阵]
	\label{g1:prop:bundle-morphism-matrices}
	设 \(E,F\) 为 Hausdorff、第二可数光滑流形 \(M\) 上的秩 \(r,s\) 的实光滑向量丛，允许 \(M\) 有边界。在共同开覆盖上取局部标架，其过渡函数为 \(g_{ij}^E,g_{ij}^F\)。一个覆盖恒等映射的逐纤维线性映射 \(\Phi:E\to F\) 为光滑丛同态，当且仅当它的局部矩阵 \(A_i:U_i\to\operatorname{Mat}_{s\times r}(\mathbb R)\) 光滑；这些矩阵满足
	\begin{equation}
		\label{g1:eq:bundle-morphism-transition}
		A_i g_{ij}^E=g_{ij}^F A_j.
	\end{equation}
	反之，任何满足此关系的光滑矩阵族都确定唯一丛同态。特别地，逐纤维为线性同构的光滑丛同态，其逆自动光滑。
\end{proposition}
\begin{proof}
	局部矩阵由 \(\Phi(e_i v)=f_i A_i v\) 定义。若 \(\Phi\) 光滑，则其作用于每个光滑基截面 \(e_{i,a}\) 后的分量光滑，故 \(A_i\) 各列光滑。反过来，光滑矩阵给出的局部映射 \((p,v)\mapsto(p,A_i(p)v)\) 光滑。交叠上，先换源分量再映射得到 \(A_i g_{ij}^E v_j\)，先映射再换目标分量得到 \(g_{ij}^F A_jv_j\)，二者对所有 \(v_j\) 相同恰好等价于\eqref{g1:eq:bundle-morphism-transition}。因此相容矩阵族唯一粘成光滑映射。最后，若每个 \(A_i(p)\) 都可逆，其逆矩阵光滑且局部逆彼此相容，从而构成光滑丛逆。
\end{proof}

例如映射的微分 \(\mathrm{T}F:\mathrm{T}M\to \mathrm{T}N\) 覆盖 \(F:M\to N\)，局部矩阵为坐标表达的导数，见\cref{g1:prop:differential-chain-rule}。底空间不同的情形只须在目标变换矩阵中代入 \(F(p)\)，矩阵判据及其证明不变。另一个常用例子是同底丛同态作用于截面：\(s\mapsto\Phi\circ s\) 把 \(\Gamma(E)\) 送入 \(\Gamma(F)\)，并与光滑函数的逐点乘法相容。

\begin{definition}
	\label{g1:def:vector-subbundle}
	设 \(E\to M\) 为秩 \(r\) 的实光滑向量丛，\(0\leq k\leq r\)。若每个 \(p\) 给定一个 \(k\) 维线性子空间 \(F_p\subset E_p\)，且每点附近存在 \(k\) 个光滑截面 \(s_1,\ldots,s_k\)，在每个邻近点都线性无关并张成相应的 \(F_p\)，则称 \(F=\coprod_p F_p\subset E\) 为秩 \(k\) 的光滑\term[zicong]{子丛}[subbundle]。
\end{definition}

定义中的光滑局部张成条件不可删去。例如在 \(\mathbb R\times\mathbb R^2\) 中，令 \(F_t=\mathbb R(1,0)\) 当 \(t\leq0\)，令 \(F_t=\mathbb R(0,1)\) 当 \(t>0\)。各个子空间的维数都为一，但若零附近有一个处处非零的连续张成截面，它在零点的值必须既落在横轴又落在纵轴，因而只能为零，矛盾。恒定维数没有替代连续变化，更没有替代光滑变化。

\begin{proposition}[子丛的适配标架]
	\label{g1:prop:subbundle-adapted-frame}
	设 \(E\to M\) 是 Hausdorff、第二可数光滑流形上的秩 \(r\) 实光滑向量丛，允许有边界。若 \(s_1,\ldots,s_k\) 在点 \(p\) 处线性无关，则在 \(p\) 的某个邻域上可将这些光滑截面补成 \(E\) 的局部标架。特别地，秩 \(k\) 子丛 \(F\subset E\) 带有唯一的光滑向量丛结构，使其局部标架给出的坐标为光滑平凡化；其拓扑是 \(E\) 的子空间拓扑，并且存在局部标架使
	\[
		E|_U\simeq U\times\mathbb R^r,\qquad
		F|_U\simeq U\times(\mathbb R^k\times\{0\}).
	\]
\end{proposition}
\begin{proof}
	先在 \(p\) 附近固定一个 \(E\) 的平凡化，把 \(s_a\) 写成光滑列向量。在 \(p\) 处用常向量 \(b_{k+1},\ldots,b_r\) 将其补为一组基，并把这些常向量在该平凡化下视为光滑截面。所成 \(r\times r\) 矩阵的行列式在 \(p\) 非零，由连续性，在某邻域仍非零。因此这些截面在整个邻域组成标架。\(k=0\) 时直接选任意标架，\(k=r\) 时无须补充。

	对子丛应用这个结论，前 \(k\) 个坐标向量正好张成 \(F\)。因此环境平凡化限制到 \(F\) 后为到 \(U\times\mathbb R^k\) 的同胚，局部拓扑正是子空间拓扑。两组这种标架的变换矩阵把 \(\mathbb R^k\times\{0\}\) 映到自身，其左上 \(k\times k\) 块为光滑可逆矩阵，故这些限制图光滑相容。Hausdorff 性与第二可数性由子空间继承，局部图给出维数 \(\dim M+k\) 及向量丛结构。唯一性由规定的局部图立即得到。边界点处仅用了光滑矩阵及非零行列式，所得局部模型仍为 \(\mathbb H^n\times\mathbb R^k\)。
\end{proof}

现在考虑一个丛同态的核与像。每条纤维内它们都是线性子空间；要成为子丛，维数先要在附近不变，而常秩条件还会给出光滑的局部基。\cref{g1:thm:constant-rank}中用过“选非零子式，再消去其余行”的办法；此处只需要对一族矩阵作消元，不需要对底空间改换常秩标准坐标。

\begin{theorem}[常秩丛同态]
	\label{g1:thm:constant-rank-bundle-morphism}
	设 \(E,F\) 为 Hausdorff、第二可数光滑流形 \(M\) 上的秩 \(r,s\) 实光滑向量丛，允许 \(M\) 有边界；\(\Phi:E\to F\) 为覆盖恒等映射的光滑丛同态。若 \(\operatorname{rank}\Phi_p=k\) 对所有 \(p\in M\) 恒成立，则
	\[
		\ker\Phi=\coprod_{p\in M}\ker\Phi_p,\qquad
		\operatorname{im}\Phi=\coprod_{p\in M}\operatorname{im}\Phi_p
	\]
	分别是 \(E,F\) 的秩 \(r-k,k\) 的光滑子丛。若秩只要求局部恒定，则此结论按各秩恒定的开集分别成立。
\end{theorem}
\begin{proof}
	固定一点并取局部标架，把 \(\Phi\) 表示为光滑 \(s\times r\) 矩阵 \(A(p)\)。先设 \(k>0\)。在基点选择一个非零 \(k\) 阶子式，重排源与目标的基后，将对应子块放到左上角。缩小邻域使该块一直可逆，写成
	\[
		A=\begin{pmatrix}B&C\\D&H\end{pmatrix},\qquad B\in\operatorname{GL}(k,\mathbb R).
	\]
	左乘可逆矩阵作消元，得到
	\[
		\begin{pmatrix}I&0\\-DB^{-1}&I\end{pmatrix}A
		=\begin{pmatrix}B&C\\0&H-DB^{-1}C\end{pmatrix}.
	\]
	再把右上块通过列消元变为零，可见右边的秩为 \(k+\operatorname{rank}(H-DB^{-1}C)\)。常秩条件遂迫使 \(H=DB^{-1}C\) 在整个邻域上成立。

	于是 \(A(x,y)=0\) 等价于 \(Bx+Cy=0\)，亦即 \(x=-B^{-1}Cy\)。所以矩阵
	\[
		\begin{pmatrix}-B^{-1}C\\I_{r-k}\end{pmatrix}
	\]
	的 \(r-k\) 个列向量给出核的光滑局部基。从下方的单位矩阵可见，这些列向量线性无关。另一方面，\(A\) 的前 \(k\) 列组成 \(\begin{psmallmatrix}B\\D\end{psmallmatrix}\)，因 \(B\) 可逆而线性无关；其余列为这些列乘以 \(B^{-1}C\)，故前 \(k\) 列正好给出像的光滑局部基。两族列向量都在先前选定标架下表示光滑截面，因而核与像满足子丛定义。

	若 \(k=0\)，矩阵处处为零，核为 \(E\)，像为零子丛。若 \(k=r\) 或 \(k=s\)，上述相应的空块与空基直接略去，结论仍成立。边界点处，\(B^{-1}\) 的矩阵元是光滑函数除以非零的 \(\det B\)，所以在半空间图中仍可光滑延拓；所有等式只在实际底空间邻域内使用，不要求延拓后的矩阵在外侧仍保持秩 \(k\)。
\end{proof}

\begin{example}
	\label{g1:exm:variable-rank-bundle-map}
	在 \(\mathbb R\) 上的两个乘积线丛之间，令 \(\Phi(t,v)=(t,tv)\)。当 \(t\ne0\) 时，它的核为零、像为整条纤维；在 \(t=0\) 时，核为整条纤维、像为零。故核与像都不能在零附近成为向量丛。即使把底空间改成 \([0,\infty)\)，同样的问题仍在边界点零发生。失败原因是秩的变化，不是边界使矩阵消元失效。
\end{example}

特别地，一个逐纤维单射的光滑丛同态把源丛同构到目标的像子丛。像是子丛由常秩定理得到；把同态看成到这个像的映射后，它逐纤维可逆，\cref{g1:prop:bundle-morphism-matrices}保证逆光滑。这允许我们把一个抽象的单射丛同态与其实现的子丛明确对应起来。
\subsection{拉回丛、直和与商丛}
\label{g1:subsec:1201:03}

设 \(f:N\to M\) 光滑，\(E\to M\) 是向量丛。如果希望在 \(N\) 的每点使用 \(f(q)\) 处的纤维，就应把新纤维定义为 \(E_{f(q)}\)。不同的 \(q\) 即使具有相同的像，也必须保留为不同的基点，因此不能只取 \(E\) 中这些纤维的并，而要把 \(q\) 一同记下来。

\begin{proposition}[拉回丛]
	\label{g1:prop:pullback-bundle}
	设 \(M,N\) 为 Hausdorff、第二可数的光滑流形，均可带边界；\(f:N\to M\) 光滑，\(\pi:E\to M\) 为秩 \(r\) 的实光滑向量丛。集合
	\[
		f^*E=\{(q,v)\in N\times E:f(q)=\pi(v)\},\qquad
		\pi'(q,v)=q
	\]
	带有自然的秩 \(r\) 实光滑向量丛结构，其拓扑是 \(N\times E\) 的子空间拓扑，称为 \(E\) 沿 \(f\) 的\term[lahuicong]{拉回丛}[pullback bundle]。映射 \(\widetilde f:f^*E\to E\)，\(\widetilde f(q,v)=v\)，是覆盖 \(f\) 的光滑丛同态。

	它具有如下泛性质：对任意 \(N\) 上的有限秩实光滑向量丛 \(H\)，以及覆盖 \(f\) 的光滑丛同态 \(\Psi:H\to E\)，存在唯一覆盖 \(\operatorname{id}_N\) 的光滑丛同态 \(\widehat\Psi:H\to f^*E\)，使 \(\widetilde f\widehat\Psi=\Psi\)。
\end{proposition}
\begin{proof}
	若 \(e_i\) 是 \(E\) 在 \(U_i\) 上的局部标架，定义
	\[
		\tau_i':f^{-1}(U_i)\times\mathbb R^r\longrightarrow(\pi')^{-1}(f^{-1}(U_i)),\qquad
		(q,a)\longmapsto(q,e_i(f(q))a).
	\]
	它是双射，逆为 \((q,v)\mapsto(q,[v]_i)\)。两个公式对所给子空间拓扑都连续，且定义域与像分别在相应局部范围内开，所以 \(\tau_i'\) 是局部平凡化的拓扑候选。它们的过渡映射为
	\[
		(q,a_j)\longmapsto(q,g_{ij}(f(q))a_j).
	\]
	这些矩阵光滑、可逆，且满足余循环条件。由\cref{g1:thm:bundle-cocycle-gluing}，得到唯一相容的光滑向量丛结构；其拓扑与上述子空间拓扑相同，因为两者在覆盖总空间的每个局部乘积上都由同一同胚给定。Hausdorff 性与第二可数性也可直接由 \(N\times E\) 的子空间性质得到。边界来自 \(N\) 的半空间坐标，纤维仍为完整的 \(\mathbb R^r\)。

	映射 \(\widetilde f\) 在这些图中写成 \((q,a)\mapsto(f(q),a)\)，故光滑，并且逐纤维线性。对泛性质，唯一可能的公式是
	\[
		\widehat\Psi(u)=\bigl(\pi_H(u),\Psi(u)\bigr).
	\]
	因为 \(\pi\Psi=f\pi_H\)，这个有序对确实落在 \(f^*E\) 中。局部矩阵中，它与 \(\Psi\) 使用相同的纤维矩阵，只是底映射变为恒等，所以光滑且逐纤维线性。此公式同时给出 \(\widetilde f\widehat\Psi=\Psi\) 与 \(\pi'\widehat\Psi=\pi_H\)，而这两个等式决定有序对的两个分量，故唯一。
\end{proof}

这两个交换关系可集中画成
\[
	\begin{tikzcd}[column sep=large,row sep=large]
		H \arrow[rr,bend left=16,"\Psi"] \arrow[dr,"\pi_H"'] \arrow[r,dashed,"\widehat\Psi"] & f^*E \arrow[r,"\widetilde f"] \arrow[d,"\pi'"'] & E \arrow[d,"\pi"] \\
		& N \arrow[r,"f"'] & M.
	\end{tikzcd}
\]
虚线箭头不增加一种任意选择；它只是把 \(\Psi(u)\) 与其新基点记录在一起。特别地，对 \(E\) 的截面 \(s\)，
\[
	(f^*s)(q)=(q,s(f(q)))
\]
是 \(f^*E\) 的光滑截面。对流形映射的微分，这一构造把覆盖 \(f\) 的 \(\mathrm{T}f:\mathrm{T}N\to \mathrm{T}M\) 改写成同底丛同态 \(\mathrm{T}N\to f^*\mathrm{T}M\)，其在 \(q\) 处仍是 \(\mathrm{d}f_q\)。纤维 \((f^*\mathrm{T}M)_q\) 是 \(\mathrm{T}_{f(q)}M\)，通常不是 \(\mathrm{T}_qN\)。

\ProseParagraphBreak

如果再有光滑 \(h:P\to N\)，则
\[
	h^*(f^*E)\longrightarrow(f\circ h)^*E,
	\qquad(p,(h(p),v))\longmapsto(p,v)
\]
是自然丛同构，其逆把 \((p,v)\) 送到 \((p,(h(p),v))\)。在从 \(E\) 拉回的局部平凡化中，两个公式的纤维分量均为恒等，故光滑。类似地，\(\operatorname{id}_M^*E\simeq E\)。这些明确同构说明，连续两次拉回与沿复合映射拉回给出相同的几何对象。

\ProseParagraphBreak

线性代数的直和则保留同一基点上的两种向量。对 \(M\) 上的秩 \(r,s\) 向量丛 \(E,F\)，定义\term[zhihecong]{直和丛}[direct sum bundle]
\[
	E\oplus F=\{(u,v)\in E\times F:\pi_E(u)=\pi_F(v)\},
	\qquad(E\oplus F)_p=E_p\oplus F_p.
\]
在共同平凡化下，它为 \(U_i\times(\mathbb R^r\oplus\mathbb R^s)\)，过渡矩阵为
\[
	\begin{pmatrix}g_{ij}^E&0\\0&g_{ij}^F\end{pmatrix}.
\]
分块乘法验证余循环条件，粘合定理给出光滑丛结构。局部表示与在 \(E\times F\) 中的子空间拓扑相容，证明与拉回丛的局部表示相同：从局部乘积到这个集合的映射及其逆都由两组平凡化给出。逐纤维投影和包含在局部都是固定线性映射，所以是光滑丛同态。这里的总空间不是整个 \(E\times F\)：直和要求两个向量属于同一基点，因而其维数为 \(\dim M+r+s\)。

\ProseParagraphBreak

商丛的目标是在每条纤维中忽略一个指定子空间的分量。虽然计算时可补选基来表示剩余分量，但商本身不应依赖所补的基，也不应要求预先选出全局补空间。

\begin{proposition}[商丛]
	\label{g1:prop:quotient-vector-bundle}
	设 \(M\) 为 Hausdorff、第二可数的光滑流形，可带边界；\(E\to M\) 为秩 \(r\) 实光滑向量丛，\(F\subset E\) 为秩 \(k\) 光滑子丛。集合
	\[
		E/F=\coprod_{p\in M}E_p/F_p
	\]
	有自然的秩 \(r-k\) 实光滑向量丛结构，其总空间拓扑为映射 \(q:E\to E/F\)，\(v_p\mapsto[v_p]\)，确定的商拓扑。\(q\) 是光滑满射丛同态，其逐纤维核为 \(F\)。如果 \(H\to M\) 是实光滑向量丛，\(\Phi:E\to H\) 是同底光滑丛同态，且 \(\Phi|_F=0\)，则存在唯一光滑丛同态 \(\overline\Phi:E/F\to H\) 使 \(\Phi=\overline\Phi q\)。
\end{proposition}
\begin{proof}
	由\cref{g1:prop:subbundle-adapted-frame}，局部可取 \(E\) 的标架 \(e_i\)，其前 \(k\) 个向量组成 \(F\) 的标架。把分量写成 \((a,b)\in\mathbb R^k\oplus\mathbb R^{r-k}\)。在每条纤维中，两个向量代表同一个商类，当且仅当它们的 \(b\) 分量相同。因此
	\[
		[e_i(p)(a,b)]\longmapsto(p,b)
	\]
	是商集合在 \(U_i\) 上到 \(U_i\times\mathbb R^{r-k}\) 的双射。

	在交叠上，两组适配标架的转换保持前 \(k\) 维子空间，故其矩阵必为
	\[
		g_{ij}=\begin{pmatrix}B_{ij}&C_{ij}\\0&D_{ij}\end{pmatrix}.
	\]
	其限制到 \(F_p\) 可逆，所以 \(B_{ij}\) 可逆；再由 \(\det g_{ij}=\det B_{ij}\det D_{ij}\ne0\)，知 \(D_{ij}\) 也可逆。商分量转换恰为 \(b_i=D_{ij}b_j\)。原余循环等式的右下块给出 \(D_{ij}D_{jk}=D_{ik}\)，于是粘合定理赋予商集合光滑向量丛结构。换用另一组适配标架时，同样的分块计算给出光滑可逆过渡，所以所得结构与补基选择无关。

	按这些局部图，\(q\) 的表达是 \((p,a,b)\mapsto(p,b)\)，故连续、光滑，并且开。开性可在每个 \(U_i\) 上验证，因为 Euclid 纤维投影把开集送到开集。具体地，对总空间开集 \(O\)，\(q(O)\) 在商丛限制到 \(U_i\) 中的部分，正是 \(O\cap E|_{U_i}\) 的局部投影像，因而开。连续开满射总是商映射：如果 \(q^{-1}(A)\) 开，则 \(A=q(q^{-1}(A))\) 开。因此这里的拓扑确实等于原集合商的商拓扑。核、满射性与秩都由同一局部公式得到。

	最后，逐纤维定义 \(\overline\Phi([v])=\Phi(v)\)。若换代表，差向量落在 \(F\)，被 \(\Phi\) 送到零，所以定义无关选择。在适配标架与 \(H\) 的任意局部标架下，\(\Phi\) 的前 \(k\) 列为零，矩阵写成 \((0,A_i)\)。\(\overline\Phi\) 的局部矩阵就是光滑矩阵 \(A_i\)，故它光滑。由于 \(q\) 满射，因子映射也唯一。
\end{proof}

对常秩丛同态 \(\Phi:E\to F\)，现在得到纤维中的第一同构定理的光滑版本：
\begin{equation}
	\label{g1:eq:bundle-first-isomorphism}
	E/\ker\Phi\ \simeq\ \operatorname{im}\Phi,
	\qquad [v_p]\longmapsto\Phi_p(v_p).
\end{equation}
核与像由\cref{g1:thm:constant-rank-bundle-morphism}成为子丛；商丛泛性质给出这个光滑映射。线性代数说明它在每条纤维上双射，而\cref{g1:prop:bundle-morphism-matrices}使其逆也光滑。因此这不是只在每个点分别存在的一族同构，而是总空间上的丛同构。

\begin{definition}
	\label{g1:def:short-exact-bundle-sequence}
	设 \(F,E,Q\) 是同一光滑流形上的有限秩实光滑向量丛，\(\iota:F\to E\)、\(p:E\to Q\) 为同底光滑丛同态。如果对每个底点 \(x\)，\(\iota_x\) 单射、\(p_x\) 满射且 \(\operatorname{im}\iota_x=\ker p_x\)，则称
	\[
		0\longrightarrow F\xrightarrow{\ \iota\ }E\xrightarrow{\ p\ }Q\longrightarrow0
	\]
	为向量丛的\term[duanzhengheconglie]{短正合列}[short exact sequence]。
\end{definition}

在这样的列中，\(\iota\) 将 \(F\) 识别为 \(E\) 的像子丛，而\eqref{g1:eq:bundle-first-isomorphism}把 \(Q\) 识别为 \(E/\iota(F)\)。局部补全 \(\iota(F)\) 的标架后，商分量可用补出的基向量提升，因而总能局部选择光滑右逆 \(Q|_U\to E|_U\)。不同邻域补出的子空间却未必彼此相同；把这些局部选择组织成整体选择，需要随后单位分解所提供的工具。即使整体选择存在，它通常也不由商丛本身唯一决定。

\subsection{标架丛与全局标架}
\label{g1:subsec:1201:04}

局部标架使向量有了分量，但一个纤维里并没有优先选定的一组基。若把所有可能的基都保留下来，换基矩阵就从计算中的辅助量变成这个新空间上的作用。这样做也能准确表达“能否在整个底空间上光滑选基”。

\begin{definition}
	\label{g1:def:frame-bundle}
	设 \(E\to M\) 为秩 \(r\) 的实光滑向量丛。定义
	\[
		\operatorname{Fr}(E)_p=\{u:\mathbb R^r\to E_p\mid u\text{ 是线性同构}\},\qquad
		\operatorname{Fr}(E)=\coprod_{p\in M}\operatorname{Fr}(E)_p.
	\]
	每个 \(u\) 等价于有序基 \((u(\varepsilon_1),\ldots,u(\varepsilon_r))\)。对任一局部标架 \(e_i\)，将 \(u\in\operatorname{Fr}(E)_p\) 唯一写成 \(u=e_i(p)A_i\)，其中 \(A_i\in\operatorname{GL}(r,\mathbb R)\)，并规定 \(u\mapsto(p,A_i)\) 为局部平凡化。若 \(e_j=e_i g_{ij}\)，则 \(A_i=g_{ij}A_j\)；用这些平凡化指定拓扑与光滑结构。所得空间称为 \(E\) 的\term[biaojiacong]{标架丛}[frame bundle]，其投影为 \(u\mapsto p\)。
\end{definition}

在 \(U_i\) 上选定参考标架 \(e_i\) 后，任意 \(u\in\operatorname{Fr}(E)_p\) 都唯一写成
\[
	u=e_i(p)A_i,\qquad A_i\in\operatorname{GL}(r,\mathbb R).
\]
因此得到 \(\operatorname{Fr}(E)|_{U_i}\simeq U_i\times\operatorname{GL}(r,\mathbb R)\)。若参考标架满足 \(e_j=e_i g_{ij}\)，则同一 \(u\) 的矩阵满足
\begin{equation}
	\label{g1:eq:frame-bundle-transition}
	A_i=g_{ij}A_j.
\end{equation}
这些过渡映射光滑可逆，并由余循环条件相容。它们赋予 \(\operatorname{Fr}(E)\) 唯一的拓扑与光滑结构，理由可直接沿用粘合证明的局部步骤：局部乘积中的开集通过过渡同胚相容，故组成拓扑基；不同基点用底空间的开邻域分离，同一基点的不同标架在同一局部乘积中分离；从底空间取可数平凡化子覆盖，再取各乘积的可数基，得到第二可数性。这里 \(\operatorname{GL}(r,\mathbb R)\) 是矩阵空间中行列式非零的开集，因此维数为 \(r^2\)，总空间维数为 \(\dim M+r^2\)。若底空间有边界，其边界仍是底空间边界的逆像。标架纤维是可逆线性映射的集合，一般不是向量空间，所以这个构造没有把 \(\operatorname{Fr}(E)\) 再当作向量丛。

\ProseParagraphBreak

群 \(\operatorname{GL}(r,\mathbb R)\) 在标架丛上有自然右作用
\begin{equation}
	\label{g1:eq:frame-bundle-right-action}
	u\cdot B=u\circ B.
\end{equation}
确为右作用，因为 \((u\cdot B)\cdot C=u\circ(BC)=u\cdot(BC)\)。局部表示为 \((p,A)\cdot B=(p,AB)\)，故光滑。作用在每条纤维内自由且传递：\(uB=u\) 蕴含 \(B=I_r\)；给定同一纤维内的 \(u,v\)，唯一使 \(uB=v\) 的矩阵为 \(B=u^{-1}\circ v\)。因此任意两组基都可通过唯一的矩阵相互转换，却没有哪一组基天然是单位元；只有选定参考标架后，纤维才被识别为矩阵群。这是以后讨论主丛时要保留的结构。

\begin{proposition}[全局标架与平凡丛]
	\label{g1:prop:global-frame-triviality}
	设 \(M\) 为 Hausdorff、第二可数光滑流形，可带边界，\(E\to M\) 为秩 \(r\) 实光滑向量丛。下列条件等价：
	\begin{equivalences}
		\item \(E\) 与乘积向量丛 \(M\times\mathbb R^r\) 光滑同构。
		\item 存在 \(r\) 个全局光滑截面，在每个点组成纤维的有序基。
		\item 标架丛投影 \(\operatorname{Fr}(E)\to M\) 有光滑截面。
	\end{equivalences}
	满足这些条件时，称 \(E\) 为\term[pingfancong]{平凡丛}[trivial bundle]。
\end{proposition}
\begin{proof}
	若给定丛同构 \(\Theta:M\times\mathbb R^r\to E\)，令 \(s_a(p)=\Theta(p,\varepsilon_a)\)。这些截面光滑，且由 \(\Theta\) 的逐纤维线性同构性构成基，故第一项推出第二项。

	给定第二项中的全局标架，定义 \(\sigma(p):\mathbb R^r\to E_p\) 为 \(\sigma(p)v=\sum_a v^a s_a(p)\)。在任一参考局部标架下，\(s_a\) 的分量组成一个光滑可逆矩阵，因此 \(\sigma\) 是标架丛的光滑截面。反过来，给定这样的 \(\sigma\)，令 \(s_a(p)=\sigma(p)\varepsilon_a\)，同一局部矩阵计算给出光滑全局标架，故第二、三项等价。

	最后，给定 \(\sigma\)，令 \(\Theta(p,v)=\sigma(p)v\)。它逐纤维为线性同构，局部表达为 \((p,v)\mapsto(p,A(p)v)\)，其中 \(A\) 光滑且可逆，所以 \(\Theta\) 与其逆 \((p,w)\mapsto(p,A(p)^{-1}w)\) 都光滑。这证明第一项。秩零时取空标架，所有构造与结论仍成立。
\end{proof}

Möbius 线丛没有处处非零的截面，因而连秩一所需的全局标架也没有。另一方面，切丛有全局标架并不要求流形具有一张全局坐标图。例如圆周上的切向量 \((-\sin t,\cos t)\) 随 \(e^{\mathrm{i}t}\) 光滑且处处非零，给出 \(\mathrm{T}S^1\) 的全局标架；圆周仍不能由一张 Euclid 开区间坐标图覆盖。标架提供的是每个点的线性坐标，流形坐标还须把邻域中的不同基点同时参数化，这两种要求应当区分。
% 纲要条目排为编号小节标题，并进入目录与书签；说明段仍保留于纲要注释中。
